
\documentclass[12pt]{article}

\usepackage{sbc-template}

\usepackage{graphicx,url}

\usepackage{orcidlink}

%\usepackage[brazil]{babel}   
\usepackage[utf8]{inputenc}  

     
\sloppy

\title{Avaliação e Mitigação de Canais Encobertos Baseados em Estados de Substituição LRU e RRIP}


\author{Adriano Araújo\inst{1}, Artur Rizzi\inst{1}, Eduardo Aniceto\inst{1}, Izabel Chaves\inst{1}\orcidlink{0009-0008-9198-3909},\\ Pedro Magalhães\inst{1}}

\address{
Departamento de Ciência da Computação -- Instituto de Ciências Exatas e Informática\\
Pontifícia Universidade Católica de Minas Gerais (PUC Minas)\\
30.535-610 -- Belo Horizonte -- Brasil\\
\email{\{iopchaves,  1464597, 1568608, 1593221, 1447017\}@pucminas.edu.br}
}

\begin{document} 

\maketitle

%-----------------------------------------------
% INTRODUÇÃO
%-----------------------------------------------
\section{Introdução}

O desempenho das arquiteturas de computadores modernas é altamente impactado pela hierarquia de memória, principalmente no nível das memórias cache, como a L1 (\textit{Level 1}), que é a mais próxima fisicamente da CPU, e também a LLC (\textit{Last-Level Cache} ou Cache de Último Nível). A utilização de caches reduz a latência de acesso à memória principal, armazenando temporariamente dados que são acessados com frequência. Entretanto, como o espaço disponível nas caches é relativamente pequeno, é necessário utilizar políticas de substituição para decidir quais blocos devem ser removidos quando um novo bloco precisa ser armazenado. Esse tipo de mecanismo também está presente em outros sistemas que possuem recursos limitados, como sistemas de gerenciamento de memória e bancos de dados, nos quais é necessário decidir quais informações devem permanecer armazenadas e quais podem ser substituídas. Entre as políticas utilizadas em caches, é destacado a LRU (\textit{Least-Recently Used}, ou Menos Recentemente Usada) e suas variações computacionalmente mais baratas, como a PLRU (\textit{Pseudo-LRU}).

No entanto, essas políticas podem introduzir vulnerabilidades de segurança, permitindo ataques de canais laterais (\textit{side-channels}) e canais encobertos (\textit{covert-channels}). Ataques como \textit{Flush+Reload} exigem uma falha de cache (\textit{cache miss}) para alterar o estado da memória. Na política LRU, porém, a própria atualização dos metadados relacionados à idade dos blocos pode vazar informações, permitindo a realização de ataques mesmo quando ocorre um acerto de cache (\textit{cache hit}). Isso possibilita que um atacante consiga inferir o padrão de acesso da vítima apenas medindo o tempo de suas próprias operações, tornando o ataque mais furtivo, mais rápido e capaz de contornar defesas baseadas no monitoramento de falhas de cache \cite{9354935}. Os processadores comerciais, a exemplo da Intel e AMD, utilizam arquiteturas de "caches seguras", que depende de particionamento ou preenchimento aleatório visando mitigar a vulnerabilidade, mas não solucionam totalmente. Como solução definitiva para essas arquiteturas, o isolamento não deve abranger apenas os dados, mas os metadados LRU também devem ser devidamente bloqueados ou particionados.

Neste contexto, a motivação desta pesquisa surge da lacuna existente na avaliação empírica e quantitativa das contramedidas propostas na literatura. Trabalhos anteriores analisaram a utilização de defesas para reduzir ataques à hierarquia de cache, enquanto estudos sobre os estados da política LRU demonstraram que esses estados podem ser explorados para estabelecer canais encobertos \cite{9354935,9065598}. No entanto, ainda falta uma avaliação rigorosa da banda residual desses canais quando mecanismos de proteção estão presentes. Além disso, ataques à LLC já demonstraram ser práticos em processadores modernos \cite{7163050}, reforçando a necessidade de investigar o comportamento desses canais em diferentes níveis da hierarquia de memória. Essa questão se torna ainda mais relevante diante do uso de políticas modernas de substituição na LLC, como a família RRIP, que utiliza estados de predição para determinar quais blocos devem ser substituídos \cite{10.1145/1815961.1815971}.

A hipótese que direciona este trabalho é que as camadas inferiores da hierarquia de memória, especificamente a LLC equipada com políticas de substituição modernas, não são totalmente imunes a ataques que exploram metadados de substituição. Além disso, isolar apenas os dados na cache é insuficiente, sendo necessário também um tratamento seguro e particionado dos estados da política para mitigar os vazamentos sem comprometer de forma significativa o desempenho.

Dessa forma, o objetivo geral deste estudo é avaliar, por meio do simulador
\textit{gem5}, o comportamento de canais encobertos baseados em políticas de
substituição de cache, comparando as políticas LRU e RRIP e analisando seu
impacto na segurança e no desempenho do sistema.

Para alcançar este propósito, são definidos os seguintes objetivos específicos:

\begin{itemize}
    \item Implementar e configurar no simulador \textit{gem5} as políticas de
    substituição LRU e RRIP para a avaliação dos canais encobertos.

    \item Implementar um canal encoberto baseado no estado da política de
    substituição e avaliar sua capacidade de transmitir informações.

    \item Comparar o comportamento das políticas LRU e RRIP utilizando métricas
    como taxa de transmissão e taxa de erro.

    \item Avaliar o impacto das políticas de substituição no desempenho do
    sistema por meio de métricas como IPC e número de acessos à cache.

    \item Avaliar uma estratégia simples de mitigação baseada no isolamento ou
    particionamento do estado de substituição e comparar seus resultados com
    as configurações sem mitigação.
\end{itemize}

\section{Revisão e Análise Crítica}

O artigo \textit{Leaking Information Through Cache LRU States in Commercial Processors and Secure Caches}\cite{9354935} investiga como os estados internos das políticas de substituição LRU podem ser explorados para criar canais laterais e canais encobertos. A principal contribuição do trabalho está em demonstrar que um acesso que resulta apenas em um \textit{cache hit} ainda é capaz de modificar o estado de substituição da cache. Dessa forma, um processo receptor pode induzir uma substituição e utilizar a diferença no tempo de acesso para inferir se outro processo acessou determinada linha de cache, possibilitando o vazamento de informações sem que o transmissor necessariamente provoque um \textit{cache miss}.

A avaliação experimental constitui um dos principais pontos fortes do artigo, pois os autores demonstram a viabilidade do ataque em processadores comerciais Intel e AMD, considerando cenários de compartilhamento por \textit{hyper-threading} e \textit{time-slicing}. Os resultados mostram taxas de transmissão significativamente maiores no cenário de \textit{hyper-threading}, chegando a aproximadamente 500 Kbps em processadores Intel, enquanto no \textit{time-slicing} as taxas são drasticamente reduzidas. Essa diferença evidencia uma limitação importante: embora o canal baseado em LRU seja viável, sua eficiência depende diretamente da arquitetura do processador e da forma como os processos compartilham os recursos da CPU.

Além da avaliação em processadores reais, o artigo analisa a vulnerabilidade de propostas de caches seguras, como a \textit{Partition-Locked Cache} e a \textit{Random Fill Cache}, utilizando o simulador gem5. Os resultados indicam que proteger ou particionar apenas os dados armazenados na cache não é suficiente quando os metadados responsáveis pela política de substituição continuam compartilhados. Entretanto, uma limitação do trabalho é que essa avaliação das caches seguras foi realizada em ambiente simulado, enquanto a validação experimental ocorreu apenas para processadores comerciais convencionais. Assim, o artigo evidencia uma lacuna relevante, que é a necessidade de investigar mecanismos de isolamento que protejam não apenas as linhas de dados, mas também os estados de substituição, avaliando de forma simultânea sua capacidade de mitigar o vazamento e seu impacto no desempenho da hierarquia de memória.

%-----------------------------------------------
% TRABALHOS RELACIONADOS
%-----------------------------------------------
\section{Trabalhos Relacionados}

\subsection{Canais e Estados de Substituição em Caches}

A pesquisa \textit{“Last-Level Cache Side-Channel Attacks are Practical”}\cite{7163050}, demonstrou o quão viável é explorar caches compartilhados, como a LLC - Last-Level Cache, em ambientes multi threaded e de nuvem para extrair chaves criptográficas. No entanto, enquanto os ataques clássicos focam no tempo do acesso à memória, o artigo em questão refina essa abordagem ao demonstrar que a própria política de substituição LRU (Least Recently Used) retém estado e é capaz do vazamento de informação. Como limitação, o ataque depende do uso de \textit{large pages} pelo atacante para montar os conjuntos de desvio (\textit{eviction sets}), e os próprios autores reconhecem que desabilitar esse recurso na VMM seria a defesa mais simples. Além disso, a arquitetura de LLC dividida em \textit{slices} obriga os autores a montar esses conjuntos por tentativa e erro. O ataque também exige etapa manual de análise e os autores admitem não ter testado o método em ambiente ruidoso ou em nuvem real.

Ademais, a análise dos componentes e estados internos do processador se conecta com \textit{“Are Coherence Protocol States Vulnerable to Information Leakage?”}\cite{8327007}, porque ambas as pesquisas compartilham a mesma ideia de que estruturas de otimização de hardware são capazes de vazarem dados através de suas alterações de estado, não se limitando ao tempo de transferência de dados. Contudo, os autores mostram que, acima de 500~Kbps, a maioria das configurações testadas perde precisão significativamente. O ataque também exige cooperação entre dois processos maliciosos, trojan e spy, o que caracteriza um canal encoberto e restringe seu uso a cenários de vazamento interno. Além disso, a criação de memória compartilhada via KSM leva dezenas de segundos, e os testes foram feitos em um único modelo de processador.

Em \textit{“Leaking Information Through Cache LRU States”}\cite{9065598} introduz-se a exploração da LRU, revelando falhas em arquiteturas que eram projetadas para isolar acessos. Porém, como qualquer acesso ao conjunto de cache altera o estado LRU, o receptor só pode medir o conjunto uma vez por rodada, obtendo no máximo um bit por medição, o que é uma limitação estrutural do canal. A política PLRU agrava isso, gerando erros de leitura. A técnica de medição por \textit{pointer chasing} também exige calibrar o tamanho da lista ligada. Por fim, a taxa de transmissão varia bastante conforme o cenário (600~Kbps em \textit{hyper-threading} contra poucos bps em compartilhamento por tempo).

Esses trabalhos adotam metodologias bem distintas entre si. \cite{7163050} avalia ataques na LLC compartilhada, implicando em modelos de ameaça onde atacar a LLC permite ataques entre VMs distintas, enquanto \cite{9065598} concentra-se no estado LRU do cache L1 privado, o que normalmente exige coexecução no mesmo núcleo. Quanto à técnica de medição, \cite{7163050} usam \textit{Prime+Probe}, \cite{9065598} usam uma variação de \textit{pointer chasing} e \cite{8327007} medem diretamente a latência de acesso a blocos em diferentes estados de coerência.

\subsection{Mitigações e Lacunas da Literatura}

No escopo das defesas, \textit{“CATalyst: Defeating last-level cache side channel attacks in cloud computing”}\cite{7446082} propõe o uso da tecnologia Cache Allocation Technology (CAT) do hardware para particionar a LLC e neutralizar ataques. O artigo dialoga com essa solução ao tratar a eficácia de mecanismos de isolamento, haja visto que as defesas puramente baseadas em software ou particionamento parcial podem ser contornadas. A limitação é que a solução depende da tecnologia Intel CAT, herdando o limite de apenas 4 \textit{Classes of Service}, insuficiente para isolar muitas VMs. O particionamento fino por página também tem limites próprios, como poucas páginas seguras por VM e perda de desempenho. Além disso, o CATalyst neutraliza ataques por contenção de linha, mas não os canais baseados em metadados de coerência ou de substituição LRU explorados em~\cite{8327007} e~\cite{9065598}.

Os ambientes experimentais das metodologias também revelam lacunas: trabalhos como \cite{7163050} e \cite{7446082} testam em máquinas virtualizadas simulando nuvem, enquanto \cite{8327007} e \cite{9065598} testam em processos nativos sem virtualização. 

Desta forma, embora Xiong et al.~\cite{9354935} tenham demonstrado a viabilidade de canais encobertos baseados no estado de substituição LRU, a avaliação de contramedidas na literatura carece de uma base comparativa quantitativa e unificada. A pesquisa original foca sua avaliação empírica via simulação (gem5) apenas nas defesas de \textit{Partition-Locked (PL) cache} e \textit{Random Fill (RF) cache}. Outras abordagens estado-da-arte, como DAWG, ScatterCache, CEASER e InvisiSpec, tiveram sua resiliência analisada de forma estritamente qualitativa. Há uma lacuna significativa na literatura quanto à medição empírica da banda residual do canal encoberto frente a essas defesas.

Além disso, estudos anteriores assumem que canais baseados em LRU em níveis mais baixos da hierarquia de memória (L2/LLC) seriam facilmente mitigáveis por técnicas preexistentes em L1. Tal suposição é feita sem uma validação experimental rigorosa e desconsidera o comportamento de políticas de substituição modernas comumente empregadas na LLC, como RRIP e suas variantes.

Para preencher essas lacunas dentro de um escopo viável, este trabalho propõe uma avaliação experimental focada, utilizando o simulador gem5 em modo \textit{syscall-emulation} (SE). Nossa metodologia adapta os protocolos de canais com e sem memória compartilhada (Algoritmos 1 e 2 de~\cite{9354935}) para medir métricas de segurança em conjunto com métricas de desempenho. 

Em suma, este artigo apresenta as seguintes contribuições principais:
\begin{itemize}
    \item \textbf{Validação Empírica na LLC:} A avaliação da capacidade do canal de substituição contra políticas modernas na LLC (como a família RRIP, nativa ou adaptada no ecossistema gem5), testando a suposição teórica da literatura de que camadas mais baixas da hierarquia seriam inerentemente mais fáceis de proteger;
    \item \textbf{Análise de Desempenho Integrada:} Uma avaliação do \textit{trade-off} entre segurança e desempenho, medindo o impacto em IPC sob cargas de trabalho abertas e controladas (ex.: PARSEC, MiBench ou microbenchmarks sintéticos);
    \item \textbf{Extensão Comparativa (Opcional):} Caso a curva de aprendizado da ferramenta permita, a implementação e medição quantitativa de uma defesa por particionamento dinâmico do estado de substituição (ex.: adaptação do conceito do DAWG) para fins de comparação direta com os dados originais.
\end{itemize}

%-----------------------------------------------
% METODOLOGIA
%-----------------------------------------------
%-----------------------------------------------
% METODOLOGIA
%-----------------------------------------------
\section{Metodologia}

Esta seção descreve o processo de construção da bancada experimental, os
protocolos de canal encoberto implementados, as variáveis e métricas
adotadas e o plano de reprodutibilidade. Todo o código e os comandos de
execução estão versionados no repositório do experimento
(\texttt{sim/}), e os resultados detalhados em \texttt{sim/results/README.md}.

\subsection{Ambiente de simulação}

Toda a avaliação foi conduzida no simulador \textit{gem5} em modo
\textit{syscall-emulation} (SE), braço X86, com a versão pinada
\texttt{f5c5a6e390f55dd5984977815bf9d0bd05da6945} (gem5 25.1, ramo
\textit{stable}). A toolchain é isolada para garantir reprodutibilidade:
Python 3.11 via \textit{mise} e SCons instalado em \texttt{sim/.venv}; o
cenário SE não requer sistema operacional convidado, o que reduz o ruído de
medição.

A hierarquia de memória simulada fixa os seguintes parâmetros:

\begin{itemize}
    \item microarquitetura: 1 núcleo \textit{TimingSimple} por defeito
          (2 contextos de hardware no cenário SMT, discutido na
          Seção~\ref{sub:canal});
    \item L1: 32 KiB, 8-way, 64 B por linha;
    \item L2: 256 KiB, 8-way;
    \item LLC: 2 MiB, 16-way;
    \item memória principal: 8 GiB;
    \item clock nominal de 2 GHz (período de 500 ps no SE).
\end{itemize}

As políticas de substituição avaliadas são: \textbf{LRU}; \textbf{SRRIP}
canônico \cite{10.1145/1815961.1815971}, obtido por parametrização da classe
\texttt{BRRIPRP} já existente no gem5 (\texttt{btp=100},
\texttt{hit\_priority=True}), sem escrever código C++ novo; \textbf{RRIP}
(\texttt{btp=100}, sem \textit{hit-priority}); e \textbf{BRRIP}
(\texttt{btp=3}, com \textit{fill-priority}). A opção
\texttt{--policy-scope} define em quais níveis a política escolhida é
aplicada (L1; L2+LLC; apenas LLC; ou todos), permitindo isolar o efeito de
cada nível na vazão do canal.

Duas restrições do simulador pinado orientaram o desenho experimental e são
documentadas como limitações da bancada:

\begin{itemize}
    \item o modelo \texttt{O3} segfaulta em
          \texttt{Decode::sortInsts} (decode.cc:488) até mesmo com um
          programa de fio único; nenhum experimento deste trabalho usa O3;
    \item no modelo \textit{TimingSimple}, a latência de uma falha
          resolvida fora da L1 não aparece no caminho
          \textit{load-to-use} (acessos que acertam L1, L2 ou DRAM medem
          tempos próximos); o único observável em medida de tempo é o par
          \textit{hit/miss} no nível L1.
\end{itemize}

\subsection{Protocolos de canal encoberto}
\label{sub:canal}

Os protocolos seguem e adaptam os algoritmos do trabalho de Xiong et al.
\cite{9354935}. Duas famílias foram implementadas em driver C para o modo
SE:

\begin{itemize}
    \item \textbf{Protocolo de fio único} (\texttt{channel.c}): um único
          processo assume os papéis de transmissor e receptor em sequência.
          A fase de \textit{decode} toca as \texttt{d} linhas do conjunto
          alvo com um ponteiro simbolicamente compartilhado
          (\textit{line[0]}); se o bit vale 1, a linha compartilhada é
          acionada uma vez. A fase de \textit{medida} percorre uma cadeia de
          ponteiros (\textit{pointer chasing}, profundidade \texttt{K})
          partindo da linha compartilhada e cronometra o percurso com
          \texttt{rdtsc}. A linha está presente na cache na fase de medida
          (\textit{hit}) sse o bit transmitido foi 1; a posição de evicção da
          política de substituição é o canal. Todos os acessos são leituras, de
          modo que o conteúdo das linhas nunca é alterado. O conjunto alvo é o
          \texttt{set 0} (stride de 4096 B entre linhas do mesmo conjunto).
    \item \textbf{Protocolo de dois fios (SMT)} (\texttt{channel\_2t.c}):
          um único binário com \texttt{pthreads} executa o receptor (fio
          principal) e o transmissor em dois \textit{thread contexts} do
          mesmo núcleo, compartilhando fisicamente a L1 (1 core,
          \texttt{numThreads=2}, CPU \textit{Timing}). O protocolo é
          \textbf{Prime+Probe}: o receptor faz o \textit{prime} de um
          conjunto de \texttt{d} linhas do anel de ponteiros; só quando o bit
          a transmitir é 1 o transmissor toca as \texttt{d} linhas do
          \textit{eviction set} (mesmo conjunto, linhas distintas), evictando
          a linha compartilhada; o receptor então cronometra o anel. A
          sincronização é feita por uma flag \texttt{volatile} compartilhada
          (\texttt{g\_phase}), sem mutex, porque o modo SE do gem5 não
          implementa \texttt{futex}. Parâmetros usados: \texttt{PRIME\_P=6}
          passes de \textit{prime}, \texttt{tail=8} rodadas descartadas e
          \texttt{warmup=32} rodadas de aquecimento.
\end{itemize}

Para executar o dois fios no SE, duas configurações de engenharia foram
necessárias: (i) \texttt{cpu.workload} precisa conter \texttt{numThreads}
processos --- reusa-se o mesmo \texttt{Process} nos dois contextos com o
override \texttt{SmtProcess} (o \texttt{initState} C++ roda uma única vez no
contexto 0; o contexto 1 nasce \textit{Halted} e é ativado pelo
chamado de sistema \texttt{clone} do \texttt{pthread\_create}); e (ii) o
binário deve ser dinâmico (\texttt{-pthread}), pois binário estático panica
com o mapeamento de TLS endereçado em página inexistente.

\subsection{Avaliação de desempenho}

Além de métricas de segurança, o trabalho mede o custo de desempenho das
políticas sob uma carga sintética de padrão \textit{scan}
(\texttt{scan.c}): um buffer do tamanho da LLC (2 MiB) é varrido em
\textit{streaming} e, entre passadas, um conjunto quente de 32 KiB
(tamanho da L1) é revisitado. Sob LRU o \textit{streaming} recicla a
L2/LLC e expulsa o conjunto quente; sob SRRIP as linhas do
\textit{streaming} entram em \textit{long distance} e o conjunto quente
conserva prioridade, reduzindo o \textit{miss rate} da LLC. Esse padrão
expõe a \textit{scan-resistance} que motiva o SRRIP e embasa a discussão do
trade-off segurança-desempenho.

\subsection{Variáveis e métricas}

\begin{itemize}
    \item \textbf{Variáveis independentes:} política de substituição
          (LRU, SRRIP, RRIP, BRRIP); escopo da política (níveis da
          hierarquia); parâmetros do canal (\texttt{d}, \texttt{K}, janela
          de medida); mensagem transmitida (fixada em 64 bits para o
          fechamento).
    \item \textbf{Variáveis dependentes:} taxa de transmissão efetiva
          (Kbps $=$ bits por rodada $\times$ 2 GHz por ciclos de rodada); taxa
          de erro medida por edit-distance Wagner-Fischer entre a mensagem
          transmitida e a recebida e por erros por bit; latências em ciclos
          (bimodalidade de \textbf{hit} vs \textbf{miss}); IPC e \textit{miss
          rate} por nível no experimento de desempenho.
\end{itemize}

A separabilidade do canal é avaliada inspecionando a distribuição bimodal de
latências e o limiar ótimo de decisão; o canal é considerado funcional quando
a edit-distance é nula nos 64 bits transmitidos.

\subsection{Reprodutibilidade}

A bancada é versionada em \texttt{sim/} com a configuração própria do gem5
(\texttt{configs/config\_channel.py}), os drivers (\texttt{src/}), os
scripts de varredura e análise (\texttt{scripts/}) e as saídas em CSV
(\texttt{results/runs/}). O build é fixado por submódulo (commit do gem5) e a
toolchain por arquivo de versão. Os comandos de execução para cada cenário ---
fio único, SMT e desempenho --- estão documentados em \texttt{sim/README.md};
todas as tabelas de resultados referenciam os scripts que as produzem.

\subsection{Uso de IA}

Nesta etapa, uma ferramenta de IA (assistente de desenvolvimento) foi
utilizada como apoio na simulação com o gem5 nas seguintes frentes, sempre
sob revisão e responsabilidade integral do grupo:

\begin{itemize}
    \item \textbf{Condução das decisões de projeto.} Uma sessão estruturada
          de entrevista (\textit{grilling}) terminou na árvore de decisões
          aprovadas pelo grupo (nomenclatura da política, cenários de
          execução, hierarquia, workloads etc.), registrada em
          \texttt{article/plano-uso-de-ia.md}.
    \item \textbf{Levantamento de fatos técnicos.} Verificação em fontes
          primárias (artigo original, código do gem5 e do ChampSim) dos
          protocolos Algoritmos 1, 2 e 3, das políticas disponíveis no gem5
          e da representação do SRRIP canônico a partir da classe
          \texttt{BRRIPRP}.
    \item \textbf{Implementação e depuração.} Escrita e validação do driver
          do canal, da configuração \textit{SE} do gem5 e dos scripts de
          métricas; em particular, a depuração das restrições do SE com
          fios (initState único via \texttt{SmtProcess},
          \texttt{workload} com \texttt{numThreads} entradas, binário
          dinâmico) e a interpretação dos desvios imprevistos (O3 quebrado,
          colapso do RRIP sob \textit{flood} cross-thread).
    \item \textbf{Análise e redação.} Geração das tabelas de resultados e
          redação desta metodologia, todo o conteúdo validado pelas métricas
          obtidas e revisado pelo grupo antes da inclusão no artigo.
\end{itemize}

As recomendações da IA acatadas incluem: representar o SRRIP canônico por
parametrização do \texttt{BRRIPRP} (dispensando código C++ novo), realizar a
medida realista no cenário SMT (L1 compartilhada) preservando
comparabilidade com o artigo original, e usar o padrão \textit{scan} para
expor a diferença de desempenho LRU $\times$ SRRIP. As hipóteses, interpretações e
a escolha dos parâmetros fixos permaneceram deliberações do grupo. O registro
completo do uso está em \texttt{article/plano-uso-de-ia.md}.

\newpage
\bibliographystyle{sbc}
\bibliography{sbc-template}

\end{document}
